\newpage
\addcontentsline{toc}{chapter}{ABSTRAK}

Afanudin, Toriq. 2026. "Pengembangan Media Pembelajaran Matematika Berbasis Simulasi 3D Menggunakan React dan Blender pada Materi Bangun Ruang Sisi Datar". Skripsi. Yogyakarta. Universitas Ahmad Dahlan

\vspace{0.5cm}

\begin{center}
    \textbf{ABSTRAK}
\end{center}

\hspace*{1cm}Penelitian ini bertujuan untuk mengembangkan media pembelajaran matematika berbasis simulasi 3D pada materi bangun ruang sisi datar serta mengetahui tingkat kelayakannya. Media pembelajaran dikembangkan menggunakan React sebagai \textit{library} pengembangan antarmuka berbasis \textit{website} dan Blender sebagai perangkat lunak untuk pembuatan objek tiga dimensi. Penelitian ini merupakan penelitian pengembangan (\textit{Research and Development}, R\&D) dengan menggunakan model pengembangan ADDIE yang terdiri atas lima tahapan, yaitu \textit{Analysis}, \textit{Design}, \textit{Development}, \textit{Implementation}, dan \textit{Evaluation}. Subjek uji coba dalam penelitian ini adalah peserta didik kelas VIII-B SMP Muhammadiyah 9 Yogyakarta. Instrumen pengumpulan data yang digunakan meliputi angket validasi ahli materi, angket validasi ahli media, dan angket respons peserta didik. Teknik analisis data yang digunakan adalah analisis deskriptif kuantitatif dengan mengkonversi data kualitatif menjadi data kuantitatif menggunakan skala Likert. Hasil penelitian menunjukkan bahwa media pembelajaran memperoleh skor rata-rata dari ahli materi sebesar \(4,63\) dengan kategori "Sangat Baik", dari ahli media sebesar \(4,66\) dengan kategori "Sangat Baik", dari uji coba kelas kecil sebesar \(4,31\) dengan kategori "Sangat Baik", dan dari uji coba kelas besar sebesar \(4,65\) dengan kategori "Sangat Baik". Rata-rata keseluruhan dari keempat penilaian tersebut adalah \(4,56\) dengan kategori "Sangat Baik". Berdasarkan hasil tersebut, dapat disimpulkan bahwa media pembelajaran matematika berbasis simulasi 3D menggunakan React dan Blender pada materi bangun ruang sisi datar dinyatakan \textbf{layak} digunakan dalam proses pembelajaran.

\vspace{0.5cm}

\textbf{Kata kunci}: media pembelajaran, simulasi 3D, React, Blender, bangun ruang sisi datar, ADDIE

\vspace{1cm}

\begin{center}
    \textbf{ABSTRACT}
\end{center}

\hspace*{1cm}This study aims to develop a 3D simulation-based mathematics learning media on the topic of polyhedron and to determine its feasibility. The learning media was developed using React as a library for website-based interface development and Blender as software for creating three-dimensional objects. This research is a Research and Development (R\&D) study using the ADDIE development model, which consists of five stages: Analysis, Design, Development, Implementation, and Evaluation. The trial subjects in this study were eighth-grade students of SMP Muhammadiyah 9 Yogyakarta. Data collection instruments included material expert validation questionnaires, media expert validation questionnaires, and student response questionnaires. The data analysis technique used was descriptive quantitative analysis by converting qualitative data into quantitative data using a Likert scale. The results showed that the learning media obtained an average score from material experts of \(4.63\) with the category "Very Good", from media experts of \(4.66\) with the category "Very Good", from the small group trial of \(4.31\) with the category "Very Good", and from the large group trial of \(4.65\) with the category "Very Good". The overall average of these four assessments was \(4.56\) with the category "Very Good". Based on these results, it can be concluded that the 3D simulation-based mathematics learning media using React and Blender on the topic of polyhedron is declared \textbf{feasible} for use in the learning process.

\vspace{0.5cm}

\textbf{Keywords}: learning media, 3D simulation, React, Blender, polyhedron, ADDIE

\newpage